\begin{keyinsight}[未来机会不等于当前EPS：Physical AI 是独立的风险调整期权]
基础估值CNY86.4、上行约3.5\%只回答截至2026年7月23日已披露汽车电子盈利的合理价值，并不穷尽德赛西威的未来机会。公司已披露AI Cube、机器人域控定点及2026年量产规划、川行致远S6低速无人车、NVIDIA Thor POC和与元戎启行的L4合作；这些事实支持建立2026--2028年的Physical AI机会树，但尚不披露客户、数量、ASP、确认收入、毛利或现金回收。正确的处理不是忽略它，也不是提前加进2026E EPS：以FY2028商业化结果按12\%折现并按55\%/35\%/10\%风险调整后，期权价值CNY2.49/股，合并参考值CNY88.9。
\end{keyinsight}

\section{两层价值框架：基础锚与未来期权必须分开}

\begin{exhibitbox}[德赛西威的两层价值]
\small
\begin{tabularx}{\textwidth}{L{2.2cm}Y Y Y}
\toprule
层次 & 可采用事实 & 当前结论 & 投资使用方式 \\
\midrule
已实现汽车智能化底座 & 2025年智能驾驶收入CNY9.700bn、同比+32.63\%，分部毛利率16.36\%；多OEM规模交付 & 已经进入合并盈利与CNY86.4基础价值 & 判断下行保护、主业质量和2026执行 \\
商业化期权：机器人与低速无人车 & AI Cube产品、机器人域控定点/规划2026量产、S6产品与未具名订单均有公司披露 & 客户、实际交付、收入、毛利和现金未披露 & 作为2026--2028事件驱动研究；基础EPS为零 \\
生态技术期权：Thor与Robotaxi & Thor方案POC、与元戎启行合作的公司公告 & 无终端项目、SOP、交付或经济学披露 & 技术能力和合作催化剂；不做收入或估值信用 \\
\bottomrule
\end{tabularx}
\sourcenote{德赛西威2025年年报、公司IR记录和公司新闻；详见本章和附录来源S11--S16。}
\end{exhibitbox}

这里的核心区别是：车端智能驾驶已经是可以验证的业务和能力底座，但不能被重复贴上“机器人收入”标签；AI Cube、低速无人车和L4生态合作证明了技术外溢的方向，却没有提供可以填入2026E利润表的财务字段。因此，3.5\%是基础业务的安全边际判断；另列的CNY2.49/股风险调整期权才是对未来机会的正式、但低置信度的估值信用。

\section{三条通往Physical AI的路径}

\subsection{路径一：高阶智驾、舱驾融合与中央计算}

这条路径最接近兑现，因为智能驾驶已形成CNY9.700bn收入。公司披露大算力、国产化芯片和跨域融合产品已批量交付，中央计算方案亦获国内多家头部客户订单。真正的未来增量不是“有订单”三个字，而是中央计算从定点/开发跨越车型SOP、持续交付、收入确认和项目毛利。高算力SoC、传感器、研发和质保也可能吞噬系统价值量；2025年智能驾驶毛利率同比下降3.55个百分点，正说明收入增长不自动转化为更高EPS。

\subsection{路径二：川行致远S6低速无人物流}

公司披露S6采用车规级开发、全栈L4系统、线控底盘及模块化上装，覆盖园区物流、城市配送和冷链运输；IR记录还提到取得相关客户订单。其未来价值在于封闭或半封闭场景中的车队部署能否把车规工程能力变成可复制交付，并形成算法数据反馈。当前仍不知道客户、车辆数量、合同是销售/租赁/运营服务、验收条件、单车价格、毛利及回款。因而市场空间和产品参数均不能替代单位经济。

\subsection{路径三：AI Cube、机器人域控制器与生态合作}

AI Cube被公司定义为复用辅助驾驶架构的机器人AI计算终端；公司IR记录确认与多家机器人企业合作、取得机器人域控项目定点，并计划于2026年量产交付。公司还披露Thor方案完成POC、与元戎启行共同开发海外L4 Robotaxi方案。这些是强于纯概念的产品和生态证据，但尚不是已完成的量产或收入。首个量产交付会提高商业化概率；只有收入、毛利和现金信息才允许它进入盈利模型。

\begin{exhibitbox}[从故事到估值的四道门]
\small
\begin{tabularx}{\textwidth}{L{2.0cm}Y Y Y}
\toprule
门槛 & 可接受的证据 & 不能替代它的材料 & 模型动作 \\
\midrule
G1 商业身份 & 具名客户/车队/项目，产品或场景及定点/合同状态明确 & ``多家合作''、产品发布或泛化订单 & 建立项目时间表，仍不给EPS \\
G2 交付转换 & SOP、客户验收、车队部署、已交付数量或首笔确认收入 & ``规划量产''、POC、样机或发布会 & 开始构建收入范围，但不将订单直接等同收入 \\
G3 单位经济 & 项目/分部毛利、价格/系统范围、增量研发及应收/库存/现金证据 & 行业TAM、技术参数或单次收入新闻 & 收入可转换为独立EPS敏感性 \\
G4 可复制性 & 多客户/多场景、复购/续约或连续期间收入与回款 & 单一试点或一次性工程项目 & 才可讨论概率加权2027--2028期权价值 \\
\bottomrule
\end{tabularx}
\sourcenote{本院基于一手披露建立的证据门槛；不是公司指引。}
\end{exhibitbox}

\section{什么规模才足以改变每股价值}

没有客户、单位、ASP、收入确认、利润率和现金流，就没有可验证的精确预测。我们先采用反向敏感性：增量EPS = 增量收入 $\times$ 净利率 $\div$ 5.968亿股；条件价值 = 增量EPS $\times$ 验证后P/E。以下数字均为说明所需规模的条件算术，\textbf{不是公司预测，也不加入2026E基础EPS}；随后在估值章节以相互排斥情景作低置信度的风险调整。

\begin{exhibitbox}[Physical AI增量价值的条件敏感性]
\small
\begin{tabularx}{\textwidth}{L{2.6cm}Y Y Y}
\toprule
年度增量收入 & 5\%净利率、20x P/E & 8\%净利率、22x P/E & 10\%净利率、25x P/E \\
\midrule
CNY0.5bn & EPS CNY0.04；CNY0.84/股 & EPS CNY0.07；CNY1.48/股 & EPS CNY0.08；CNY2.10/股 \\
CNY1.0bn & EPS CNY0.08；CNY1.68/股 & EPS CNY0.13；CNY2.95/股 & EPS CNY0.17；CNY4.19/股 \\
CNY2.0bn & EPS CNY0.17；CNY3.35/股 & EPS CNY0.27；CNY5.90/股 & EPS CNY0.34；CNY8.38/股 \\
CNY5.0bn & EPS CNY0.42；CNY8.38/股 & EPS CNY0.67；CNY14.74/股 & EPS CNY0.84；CNY20.94/股 \\
\bottomrule
\end{tabularx}
\sourcenote{本院条件敏感性，股本为5.968亿股；P/E为商业化证据充分后的情景倍数，不代表当前市场对新业务的定价。}
\end{exhibitbox}

敏感性给出一个务实结论：CNY0.5--1.0bn的新增业务在战略上重要，但在硬件型利润率下未必足以使CNY49.82bn市值出现大幅重估；真正具有凸性的结果需要数十亿元级的可重复收入、可验证利润和更长的增长持续期。这是2027--2028年需要前瞻研究的赔率，不是2026年应当预支的利润。为使它成为可比较的当前价值，我们仅对FY2028两种商业结果给出35\%和10\%权重，另以55\%对应未形成可衡量商业转化；其概率加权现值为CNY2.49/股。

\section{未来机会的行动框架}

Physical AI使德赛西威从单一的``市场支持型观察''增加了一条\stance{riskamber}{未来机会催化剂观察线}。最有信息量的节点是：机器人域控2026年规划量产是否真实发生并出现验收/收入；S6是否披露付费场景、车队交付、商业模式和回款；中央计算是否进入SOP/收入且智驾毛利止跌；新业务扩张是否没有以应收、库存和经营现金流恶化为代价。若没有这些证据，即使主题热度上升，也应维持零基础EPS并下调风险调整权重；若连续跨过G1--G3，则应重建独立的2027--2028情景，而非仅把主业P/E上调。

\begin{exhibitbox}[未来机会的证据记分卡]
\small
\begin{tabularx}{\textwidth}{L{2.1cm}Y Y}
\toprule
观察项 & 事实升级的最低标准 & 投资含义 \\
\midrule
机器人域控 & 2026年计划量产出现实际交付/验收；继而披露客户、收入或项目毛利 & 从技术期权转为商业化跟踪；收入前仍不给基础EPS \\
川行致远S6 & 付费场景、车队交付与合同模式可核验；再观察复购和回款 & 判定是产品展示、一次性项目还是可复制的交付/服务业务 \\
中央计算/舱驾融合 & 匿名获单进入车型SOP或确认收入，且智驾毛利率不再下行 & 可以提高汽车智能化成长桥可信度；先不以订单额替代收入 \\
现金纪律 & 新业务推进期间应收、库存、减值和CFO不显著劣化 & 防止以现金占用换取表面增长；若恶化则折减期权价值 \\
\bottomrule
\end{tabularx}
\sourcenote{本院监测框架；升级需依赖公司、客户或交易所可审计资料。}
\end{exhibitbox}

\begin{riskbox}[研究结论的边界]
Physical AI应当在利润表兑现前被持续研究，因为首个可核验的商业化事实可能改变2027--2028年的预期；但它不应替代主业毛利、现金转换和估值安全边际的纪律。若2026年后仍未出现交付、收入或单位经济证据，最合理的更新是降低期权权重，而非以更远期叙事延长估值期限。
\end{riskbox}
